%----------------------------------------------------------------------------------------
% VIERKANTSWORTELS EN VERGELIJKINGEN. Alle opgaven zijn eigen werk.
%
% Opbouw: Start, leerstof (theorie, werkwijze of uitleg, voorbeeld, korte oefening), en op een
% nieuwe pagina de Opdrachten (door elkaar, stijgende moeilijkheid), Koppel en Test jezelf.
% Labels van dit hoofdstuk beginnen met h2-.
%----------------------------------------------------------------------------------------

\mychapter{2}{Vierkantswortels en vergelijkingen}

\lkr{Leerplan B+S'': LPD 40 (tweedegraadsvergelijkingen met reële coëfficiënten oplossen in $\C$), met de wenk bij LPD 39 (vierkantswortels van een complex getal in cartesische vorm). De vergelijkingen van hogere graad in sectie 3 lossen leerlingen alleen op door af te zonderen en te ontbinden, met de technieken uit de 2de graad. Het wortelteken staat nooit bij een negatief getal.}

In hoofdstuk 1 bleek dat $2 + 3i$ een oplossing is van $x^2 - 4x + 13 = 0$, een vergelijking zonder reële oplossingen. Dat was geen toeval. In dit hoofdstuk zie je dat elke tweedegraadsvergelijking met reële coëfficiënten oplossingen heeft in $\C$, en hoe je ze vindt: met dezelfde discriminant als vroeger. De oplossingen komen in paren van toegevoegde getallen. Op het einde zoek je ook de vierkantswortels van een complex getal zoals $5 + 12i$.

\begin{oefeningen}[Start]
\oef{h2-start-1}

\oef{h2-verken}
\end{oefeningen}

\lkr{Laat oefening \ref{oef:h2-verken} maken vóór de uitleg. In vraag c lossen leerlingen een tweedegraadsvergelijking op zonder formule, door een kwadraat af te splitsen: dat is het bewijs van de stelling in sectie 2, met getallen. Vraag d bereidt som en product voor.}

%========================================================================================
\section{Vierkantswortels van een negatief getal}

In $\R$ heeft een negatief getal geen vierkantswortels. In $\C$ wel: $3i$ en $-3i$ hebben allebei kwadraat $-9$. We breiden het begrip vierkantswortel daarom uit naar complexe getallen.

\begin{theorie}
Een complex getal $w$ is een \textbf{\textcolor{accent}{vierkantswortel}} van $z$ als \quad $\formule{w^2 = z}$

\begin{stelling}{Vierkantswortels van een negatief getal}
Voor een reëel getal $r > 0$ heeft $-r$ twee vierkantswortels: \quad $\formule{i\sqrt{r}}$ \ en \ $\formule{-i\sqrt{r}}$
\end{stelling}

\begin{bewijs}
$(\pm i\sqrt{r})^2 = i^2 \cdot \bigl(\sqrt{r}\bigr)^2 = -r$. Er zijn er niet meer: als $w^2 = -r$, dan is
\bewijsregel{w^2 - \bigl(i\sqrt{r}\bigr)^2 = 0, \quad \text{dus} \quad \bigl(w - i\sqrt{r}\bigr)\bigl(w + i\sqrt{r}\bigr) = 0 ,}
en een product is alleen 0 als een van de factoren 0 is, dus $w = i\sqrt{r}$ of $w = -i\sqrt{r}$.
\end{bewijs}
\end{theorie}

Het bewijs gebruikt dat een product van complexe getallen alleen 0 is als een factor 0 is. Dat geldt in $\C$ net als in $\R$: is $zw = 0$ en $z \neq 0$, dan vermenigvuldig je met $\frac{1}{z}$ en vind je $w = 0$.

We schrijven de vierkantswortels als $\pm i\sqrt{5}$ en nooit als $\sqrt{-5}$: zoals je in hoofdstuk 1 zag, leidt het wortelteken bij een negatief getal tot fouten. Je schrijft $i$ liefst vóór de wortel, zodat je niet leest dat $i$ onder het wortelteken staat.

\voorbeeld{}{Bepaal de vierkantswortels van $-12$.\label{vb:h2-wortel}}{%
\uitleg{Het getal is negatief, dus de vierkantswortels zijn $\pm i\sqrt{12}$. We vereenvoudigen de wortel.}
\begin{stappen}
\pm i\sqrt{12} = \pm i\sqrt{4 \cdot 3} = \pm 2i\sqrt3 & $\sqrt{12} = 2\sqrt3$. \\[1ex]
\bigl(2i\sqrt3\bigr)^2 = 4 \cdot i^2 \cdot 3 = -12 & Controle.
\end{stappen}}

\begin{oefeningen}
\oef{h2-wortel}
\end{oefeningen}

%========================================================================================
\section{Tweedegraadsvergelijkingen}

\subsection{De discriminant}

Een tweedegraadsvergelijking $az^2 + bz + c = 0$ los je op met de discriminant $D = b^2 - 4ac$. Was $D < 0$, dan was er in $\R$ geen oplossing, omdat je de vierkantswortel van $D$ nodig had. In $\C$ heeft $D$ wel vierkantswortels.

\begin{theorie}
\begin{stelling}{Tweedegraadsvergelijking met reële coëfficiënten}
De vergelijking $az^2 + bz + c = 0$ met $a, b, c \in \R$, $a \neq 0$ en discriminant $D = b^2 - 4ac$ heeft in $\C$

\vspace{1mm}
als $D \geqslant 0$: de reële oplossingen \quad $\formule{z = \dfrac{-b \pm \sqrt{D}}{2a}}$

\vspace{1mm}
als $D < 0$: de twee toegevoegde oplossingen \quad $\formule{z = \dfrac{-b \pm i\sqrt{-D}}{2a}}$
\end{stelling}

\begin{bewijs}
Deel door $a$ en splits een kwadraat af, zoals in de Verken:
\bewijsregel{z^2 + \frac{b}{a}z + \frac{c}{a} = 0 \iff \Bigl(z + \frac{b}{2a}\Bigr)^2 - \frac{b^2}{4a^2} + \frac{c}{a} = 0 \iff \Bigl(z + \frac{b}{2a}\Bigr)^2 = \frac{D}{4a^2} .}
Is $D < 0$, dan is $\frac{D}{4a^2}$ negatief, met vierkantswortels $\pm i\sqrt{\frac{-D}{4a^2}} = \pm \frac{i\sqrt{-D}}{2|a|}$. Dus $z + \frac{b}{2a} = \pm\frac{i\sqrt{-D}}{2a}$ (het $\pm$-teken vangt het teken van $a$ op), en dat geeft de formule. Is $D \geqslant 0$, dan vind je zo de formule die je al kent.
\end{bewijs}
\end{theorie}

Je hoeft dus maar één formule te kennen: als $D < 0$ vervang je $\sqrt{D}$ door $i\sqrt{-D}$. In $\C$ heeft een tweedegraadsvergelijking met reële coëfficiënten altijd oplossingen: twee als $D \neq 0$, één als $D = 0$.

\voorbeeld{}{Los op in $\C$: $z^2 - 4z + 13 = 0$.\label{vb:h2-vkv}}{%
\uitleg{We berekenen de discriminant. Die is negatief, dus vervangen we $\sqrt{D}$ door $i\sqrt{-D}$.}
\begin{stappen}
D = (-4)^2 - 4 \cdot 1 \cdot 13 = 16 - 52 = -36 & $D < 0$: twee toegevoegde oplossingen. \\[1ex]
z = \frac{4 \pm i\sqrt{36}}{2} = \frac{4 \pm 6i}{2} & De vierkantswortels van $-36$ zijn $\pm 6i$. \\[3ex]
z = 2 + 3i \quad \text{of} \quad z = 2 - 3i & Deel reëel en imaginair deel door 2.
\end{stappen}}

\voorbeeld{}{Los op in $\C$: $2z^2 + 2z + 5 = 0$.\label{vb:h2-vkv-a}}{%
\uitleg{Nu is $a = 2$: let op de noemer $2a = 4$.}
\begin{stappen}
D = 2^2 - 4 \cdot 2 \cdot 5 = 4 - 40 = -36 & Weer $D < 0$. \\[1ex]
z = \frac{-2 \pm 6i}{4} = -\frac12 \pm \frac32 i & Splits in reëel en imaginair deel.
\end{stappen}}

\begin{werkwijze}
Breng alles naar één lid en bereken $D = b^2 - 4ac$. Is $D < 0$, zoek dan de vierkantswortels $\pm i\sqrt{-D}$ en vul ze in in de formule, in de plaats van $\pm\sqrt{D}$. Schrijf de oplossingen in de vorm $a + bi$. Is $b = 0$, dan zonder je $z^2$ af: $z^2 + 16 = 0$ geeft $z^2 = -16$, dus $z = \pm 4i$.
\end{werkwijze}

\nujij{Los op in $\C$: $z^2 + 2z + 2 = 0$.}{$D = 4 - 8 = -4$, dus $z = \frac{-2 \pm 2i}{2} = -1 \pm i$.}

\naast{\begin{minipage}{10cm}
Waar zitten die oplossingen in een grafiek? De parabool met vergelijking $y = x^2 - 4x + 13$ snijdt de $x$-as niet: daarom zijn er geen reële oplossingen. De complexe oplossingen $2 \pm 3i$ zie je niet als snijpunten, maar hun reëel deel 2 is wel de $x$-coördinaat van de top. In Koppel ontdek je wat het imaginair deel met de grafiek te maken heeft.
\end{minipage}}{\begin{tikzpicture}[x=.5cm,y=.2cm,font=\scriptsize]
  \draw[black!14,very thin,xstep=1,ystep=5] (-1,0) grid (5,20);
  \draw[->,grijs] (-1,0) -- (5.5,0) node[right,text=black] {$x$};
  \draw[->,grijs] (0,0) -- (0,21.5) node[above,text=black] {$y$};
  \foreach \t in {2,4}{\node[below,inner sep=1.5pt] at (\t,0) {$\t$};}
  \foreach \t in {5,10,15,20}{\node[left,inner sep=1.5pt] at (0,\t) {$\t$};}
  \draw[accent,thick] plot[domain=-.6:4.6,samples=60,smooth] (\x,{(\x)^2-4*(\x)+13});
  \fill[accent] (2,9) circle (2pt) node[below,text=black,inner sep=3pt] {top $(2, 9)$};
\end{tikzpicture}}

\begin{oefeningen}
\oef{h2-vkv}
\end{oefeningen}

%========================================================================================
\section{Toegevoegde oplossingen}

\subsection{Rekenregels voor de toegevoegde}

De twee niet-reële oplossingen van een tweedegraadsvergelijking met reële coëfficiënten zijn altijd elkaars toegevoegde. Dat volgt uit de formule, maar het geldt ook voor vergelijkingen van hogere graad. Daarvoor hebben we eerst rekenregels voor de toegevoegde nodig.

\begin{theorie}
\begin{stelling}{Rekenregels voor de toegevoegde}
Voor complexe getallen $z$ en $w$, een reëel getal $r$ en een natuurlijk getal $n$:

\vspace{1mm}
\qquad $\formule{\conj{z + w} = \conj{z} + \conj{w}} \qquad \formule{\conj{z \cdot w} = \conj{z} \cdot \conj{w}} \qquad \formule{\conj{z^n} = \conj{z}^{\,n}} \qquad \formule{\conj{r} = r}$
\end{stelling}

\begin{bewijs}
Neem $z = a + bi$ en $w = c + di$. Voor het product is $zw = (ac - bd) + (ad + bc)i$, dus
\bewijsregel{\conj{z \cdot w} = (ac - bd) - (ad + bc)i \quad \text{en} \quad \conj{z} \cdot \conj{w} = (a - bi)(c - di) = (ac - bd) - (ad + bc)i .}
De som gaat op dezelfde manier. De regel voor $z^n$ volgt uit die voor het product, met $n$ gelijke factoren. Een reëel getal heeft imaginair deel 0, dus $\conj{r} = r$.
\end{bewijs}

\begin{stelling}{Toegevoegde oplossingen}
Is $z$ een oplossing van een veeltermvergelijking met reële coëfficiënten, dan is ook $\conj{z}$ een oplossing.
\end{stelling}

\begin{bewijs}
We doen het voor graad 3; voor een andere graad gaat het net zo. Stel $az^3 + bz^2 + cz + d = 0$ met $a, b, c, d \in \R$. Neem van beide leden de toegevoegde, en gebruik de rekenregels:
\bewijsregel{\conj{a}\,\conj{z}^{\,3} + \conj{b}\,\conj{z}^{\,2} + \conj{c}\,\conj{z} + \conj{d} = \conj{0} \quad \text{dus} \quad a\conj{z}^{\,3} + b\conj{z}^{\,2} + c\conj{z} + d = 0 ,}
want de coëfficiënten en 0 zijn reëel. Dus $\conj{z}$ is ook een oplossing.
\end{bewijs}
\end{theorie}

Ken je dus één niet-reële oplossing, dan ken je er twee. Let op: de coëfficiënten moeten reëel zijn. In opdracht \ref{oef:h2-o15} zie je een vergelijking met een niet-reële coëfficiënt, waarvan de oplossingen geen toegevoegden zijn.

\subsection{Som, product en ontbinden}

In de Verken was de som van de oplossingen $-3 + 2i$ en $-3 - 2i$ gelijk aan $-6$, en hun product $13$: de coëfficiënten van $z^2 + 6z + 13$, op het teken na. Dat ken je al bij reële oplossingen, en het blijft gelden in $\C$.

\begin{theorie}
\begin{stelling}{Som, product en ontbinding}
Heeft $az^2 + bz + c = 0$ de oplossingen $z_1$ en $z_2$ (eventueel gelijk), dan is

\vspace{1mm}
\qquad $\formule{z_1 + z_2 = -\dfrac{b}{a}} \qquad \formule{z_1 \cdot z_2 = \dfrac{c}{a}} \qquad \formule{az^2 + bz + c = a(z - z_1)(z - z_2)}$
\end{stelling}

\begin{bewijs}
Uit de formule is $z_1 + z_2 = \frac{-b + \delta}{2a} + \frac{-b - \delta}{2a} = -\frac{b}{a}$, met $\delta = \sqrt{D}$ of $\delta = i\sqrt{-D}$, en in beide gevallen $\delta^2 = D$. Dan is
\bewijsregel{z_1 z_2 = \frac{(-b + \delta)(-b - \delta)}{4a^2} = \frac{b^2 - \delta^2}{4a^2} = \frac{b^2 - b^2 + 4ac}{4a^2} = \frac{c}{a} .}
Werk je $a(z - z_1)(z - z_2) = a\bigl(z^2 - (z_1 + z_2)z + z_1 z_2\bigr)$ uit, dan vind je zo $az^2 + bz + c$.
\end{bewijs}
\end{theorie}

Bij toegevoegde oplossingen $p \pm qi$ is de som $2p$ en het product $p^2 + q^2$, de modulus in het kwadraat. Zo stel je snel een vergelijking op.

\voorbeeld{}{Stel een tweedegraadsvergelijking met reële coëfficiënten op die $1 - 4i$ als oplossing heeft.\label{vb:h2-opstellen}}{%
\uitleg{De coëfficiënten zijn reëel, dus de andere oplossing is de toegevoegde. Met de som en het product kennen we de vergelijking.}
\begin{stappen}
z_2 = 1 + 4i & De toegevoegde oplossing. \\[1ex]
z_1 + z_2 = 2, \quad z_1 z_2 = 1^2 + 4^2 = 17 & Som en product. \\[1ex]
z^2 - 2z + 17 = 0 & Met $a = 1$: $b = -2$ en $c = 17$.
\end{stappen}
Elk veelvoud, zoals $3z^2 - 6z + 51 = 0$, is ook goed.}

\voorbeeld{}{Ontbind $z^2 + 2z + 10$ in factoren van de eerste graad.\label{vb:h2-ontbinden}}{%
\uitleg{In $\R$ lukt dat niet, want de discriminant is negatief. In $\C$ zoeken we de oplossingen van $z^2 + 2z + 10 = 0$, en gebruiken we de ontbinding uit het kader.}
\begin{stappen}
D = 4 - 40 = -36, \quad z = \frac{-2 \pm 6i}{2} = -1 \pm 3i & De oplossingen. \\[3ex]
z^2 + 2z + 10 = (z + 1 - 3i)(z + 1 + 3i) & $z - (-1 + 3i) = z + 1 - 3i$.
\end{stappen}
Controle: $(z + 1)^2 - (3i)^2 = z^2 + 2z + 1 + 9$.}

\nujij{Stel een vergelijking op met oplossing $5i$.}{De andere oplossing is $-5i$; som 0, product 25: $z^2 + 25 = 0$.}

\begin{oefeningen}
\oef{h2-opstellen}
\end{oefeningen}

\subsection{Vergelijkingen van hogere graad}

Een vergelijking van hogere graad los je op zoals in $\R$: je brengt ze terug tot vergelijkingen van lagere graad, door een factor af te zonderen of te ontbinden. Het verschil is dat een factor als $z^2 + 4$ nu ook oplossingen geeft.

\voorbeeld{}{Los op in $\C$: $z^3 + z^2 + 4z + 4 = 0$.\label{vb:h2-hoger}}{%
\uitleg{We ontbinden door termen te groeperen: de eerste twee termen hebben $z^2$ gemeen, de laatste twee 4.}
\begin{stappen}
z^2(z + 1) + 4(z + 1) = 0 & Zonder in de groepen een factor af. \\[1ex]
(z + 1)(z^2 + 4) = 0 & De factor $z + 1$ komt twee keer voor. \\[1ex]
z = -1 \quad \text{of} \quad z^2 = -4 & Een product is 0 als een factor 0 is. \\[1ex]
z = -1 \quad \text{of} \quad z = 2i \quad \text{of} \quad z = -2i & Drie oplossingen, waarvan twee toegevoegd.
\end{stappen}}

\begin{werkwijze}
Ontbind de veelterm: zonder een factor af, groepeer termen, gebruik een merkwaardig product, of deel door $z - r$ als je een oplossing $r$ kent. Los daarna elke factor op, ook de tweedegraadsfactoren met $D < 0$.
\end{werkwijze}

\begin{oefeningen}
\oef{h2-hoger}
\end{oefeningen}

%========================================================================================
\section{Vierkantswortels van een complex getal}

Elk negatief getal heeft twee vierkantswortels in $\C$. Hoe zit het met een niet-reëel getal, zoals $5 + 12i$? In de Opwarmer heb je $(3 + 2i)^2 = 5 + 12i$ uitgerekend, dus $3 + 2i$ is een vierkantswortel van $5 + 12i$, en $-(3 + 2i)$ ook. Maar hoe vind je die zonder ze te raden?

\begin{theorie}
\begin{stelling}{Vierkantswortels}
Elk complex getal $z \neq 0$ heeft precies twee vierkantswortels $\formule{w}$ en $\formule{-w}$.
\end{stelling}
\end{theorie}

Dat er hoogstens twee zijn, bewijs je zoals in sectie 1: is $v^2 = w^2$, dan is $(v - w)(v + w) = 0$, dus $v = w$ of $v = -w$. Dat er altijd een vierkantswortel bestaat, zie je in de werkwijze na de voorbeelden: ze werkt voor elk getal, want $x^2 = \frac12\bigl(|z| + \Re z\bigr)$ en $y^2 = \frac12\bigl(|z| - \Re z\bigr)$ zijn nooit negatief. In hoofdstuk 4 vind je de vierkantswortels ook meetkundig.

Stel $w = x + yi$ met $x, y \in \R$. Dan is $w^2 = (x^2 - y^2) + 2xyi$. Uit $w^2 = z$ volgen twee vergelijkingen. Een derde vergelijking haal je uit de modulus: $|w|^2 = |w^2| = |z|$ (zie opdracht \ref{oef:h1-o18} van hoofdstuk 1), dus $x^2 + y^2 = |z|$.

\voorbeeld{}{Bepaal de vierkantswortels van $5 + 12i$.\label{vb:h2-vw}}{%
\uitleg{We zoeken $w = x + yi$ met $w^2 = 5 + 12i$. Uit het reëel deel, het imaginair deel en de modulus halen we drie vergelijkingen voor $x$ en $y$.}
\begin{stappen}
x^2 - y^2 = 5, \quad 2xy = 12 & Reëel en imaginair deel van $w^2 = (x^2 - y^2) + 2xyi$. \\[1ex]
x^2 + y^2 = \sqrt{5^2 + 12^2} = 13 & De modulus: $|w|^2 = |5 + 12i|$. \\[1ex]
2x^2 = 18, \quad 2y^2 = 8 & Tel de eerste en de derde vergelijking op, en trek ze af. \\[1ex]
x = \pm 3, \quad y = \pm 2 & Nog niet elke combinatie is goed. \\[1ex]
w = 3 + 2i \quad \text{of} \quad w = -3 - 2i & $xy = 6 > 0$: $x$ en $y$ hebben hetzelfde teken.
\end{stappen}}

\voorbeeld{}{Bepaal de vierkantswortels van $-8 - 6i$.\label{vb:h2-vw2}}{%
\uitleg{We volgen dezelfde stappen. Nu is $2xy$ negatief, en dat bepaalt welke tekens bij elkaar horen.}
\begin{stappen}
x^2 - y^2 = -8, \quad 2xy = -6, \quad x^2 + y^2 = 10 & Want $|-8 - 6i| = \sqrt{64 + 36} = 10$. \\[1ex]
x^2 = 1, \quad y^2 = 9 & Optellen en aftrekken. \\[1ex]
w = 1 - 3i \quad \text{of} \quad w = -1 + 3i & $xy = -3 < 0$: verschillend teken.
\end{stappen}
Controle: $(1 - 3i)^2 = 1 - 6i + 9i^2 = -8 - 6i$.}

\begin{werkwijze}
Stel $w = x + yi$. Schrijf op: $x^2 - y^2 = \Re z$, \ $x^2 + y^2 = |z|$ \ en \ $2xy = \Im z$. Uit de eerste twee vind je $x^2$ en $y^2$. Het teken van $xy$ zegt welke tekens van $x$ en $y$ samenhoren. Controleer door te kwadrateren.
\end{werkwijze}

\nujij{Bepaal de vierkantswortels van $3 - 4i$.}{$x^2 - y^2 = 3$, $x^2 + y^2 = 5$, $xy = -2$: $x^2 = 4$, $y^2 = 1$, dus $w = \pm(2 - i)$.}

\begin{oefeningen}
\oef{h2-vw}
\end{oefeningen}

%========================================================================================
\opdrachtenkop

\begin{oefeningen}
\oef{h2-o1}

\oef{h2-o2}

\oef{h2-o3}

\oef{h2-o4}

\oef{h2-o5}

\oef{h2-o6}

\oef{h2-o7}

\oef{h2-o8}

\oef{h2-o9}

\oef{h2-o10}

\oef{h2-o11}

\oef{h2-o12}

\oef{h2-o13}

\oef{h2-o14}

\oef{h2-o15}

\oef{h2-o16}

\oef{h2-o17}

\oef{h2-o18}
\end{oefeningen}

%----------------------------------------------------------------------------------------
% KOPPEL: verbanden met andere leerstof
%----------------------------------------------------------------------------------------
\needspace{20\baselineskip}
\opdrachtsectie{Koppel}

\lkr{Tweedegraadsfuncties kennen de leerlingen uit de 2de graad. De tweede oefening is een stukje geschiedenis: zo kwamen de complexe getallen in de wiskunde.}

\begin{oefeningen}
\oef{h2-koppel-1}

\oef{h2-koppel-2}
\end{oefeningen}

%----------------------------------------------------------------------------------------
\opdrachtsectie{Test jezelf}

\begin{oefeningen}
\oef{h2-t1}

\oef{h2-t2}

\oef{h2-t3}

\oef{h2-t4}

\oef{h2-t5}

\oef{h2-t6}

\oef{h2-t7}
\end{oefeningen}

\herhaalschema{%
  \fv{oef:h2-t1}{vb:h2-wortel}{oef:h2-wortel}%
  \fv{oef:h2-t2}{vb:h2-vkv}{oef:h2-vkv}%
  \fv{oef:h2-t3}{vb:h2-vkv-a}{oef:h2-vkv}%
  \fv{oef:h2-t4}{vb:h2-opstellen}{oef:h2-opstellen}%
  \fv{oef:h2-t5}{vb:h2-ontbinden}{oef:h2-opstellen}%
  \fv{oef:h2-t6}{vb:h2-hoger}{oef:h2-hoger}%
  \fv{oef:h2-t7}{vb:h2-vw}{oef:h2-vw}}

\ifzelfstudie
\section*{Maak je eigen samenvatting}
Sluit dit boek en schrijf in je schrift: wat een vierkantswortel van een complex getal is, de vierkantswortels van een negatief getal, hoe je een tweedegraadsvergelijking oplost als de discriminant negatief is, de rekenregels voor de toegevoegde, waarom niet-reële oplossingen in toegevoegde paren komen, som, product en ontbinding, hoe je een vergelijking van hogere graad aanpakt, en hoe je de vierkantswortels van een complex getal vindt. Vergelijk daarna met de kaders.
\fi
